\newpage
\section{Introduction} \label{sec:introduction}

%=========================================================

% SPACE ---   \hspace{0.6em}   ---
% GLOSSARY ---   \gls{gt}   ---
% BOLD ---   \textbf{mathematics}   ---

%=========================================================

\subsection{Motivation and Problem Statement}

\hspace{0.6em} As a \gls{swe}, software quality, development efficiency and productivity are important aspects of my day-to-day work. Writing and maintaining software also requires identifying defects, reviewing existing code and ensuring that changes do not introduce new problems. These activities can be time consuming, especially in larger projects. This has motivated me to explore how \gls{ai}, Machine Learning (\gls{ml}) and \gls{dl} can support software engineering activities and improve the way developers work.

\hspace{0.6em} One of the various areas these technologies have significant potential is \textbf{static bug detection}. \textit{Testing can reveal failures by executing a program with particular inputs}, \textbf{static analysis examines source code or compiled code without executing the program} \cite{paperreference1_ayewah2008using}. This allows potential defects to be identified during development and makes \textbf{static analysis useful in code reviews, development environments and automated build processes}.

\hspace{0.6em} Traditional static analysis also has limitations. Static analyzers can identify potential problems but also produce warnings that do not represent actual defects. These are known as \emph{false positives}. Developers must manually investigate such warnings, which consumes time and can reduce trust in the analysis tools. \textbf{\gls{ml} offers an opportunity to improve this process by learning patterns from source code, historical defects and previously analyzed warnings}. These can complement existing static analysis tools rather than completely replace them.

\hspace{0.6em} The purpose of this research is to explore how \gls{ml} techniques, including \gls{dl} and more recent \gls{ai} based approaches, can improve static bug detection and support software developers in their daily activities. 


%=========================================================

\subsection{Research Scope}

\hspace{0.6em} This report examines traditional static bug detection and the application of \gls{ml} techniques to reduce false positives and improve warning classification.

\hspace{0.6em} The theoretical foundation is established through \textit{Using Static Analysis to Find Bugs} by \textit{Ayewah et al.}, which describes the \textbf{FindBugs} approach and its practical use of bytecode inspection, data flow analysis and bug pattern detection \cite{paperreference1_ayewah2008using}. The industrial study in paper "\textit{A few billion lines of code later...}" by \textit{Bessey et al.} provides additional context regarding scalability, false positives and developer trust \cite{paperreference2_bessey2010few}.

\hspace{0.6em} The \gls{ml} analysis focuses on three complementary studies. The first study investigate a hybrid approach that uses feature and Transformer based models to filter false positive warnings generated by the \textbf{Infer} \cite{meta_infer} static analyzer \cite{paper1_kharkarlearning2022}. The second study propose a large real world dataset of \textit{SonarQube} \cite{sonarqube_docs} warnings and evaluate lightweight code embeddings with \gls{ml} classifiers for false positive filtering \cite{paper2_hegedHus2022static}. The last paper compare traditional \gls{ml} models, recurrent \gls{nns} and graph neural networks (\gls{ggnns}) for classifying static analysis reports using different source code representations 
\cite{paper3_yerramreddy2023empirical}. Together, these studies provide complementary perspectives on warning classification, source code representation, training data and the effectiveness of \gls{ml} assisted static analysis.
